\chapter{Implémentation}
\label{ch:impl}

\section{Backend FastAPI}
Le backend expose une API REST documentée automatiquement via
OpenAPI à l'URL \texttt{/docs}. \Cref{tab:endpoints} résume les
principaux endpoints.

\begin{table}[H]\centering\small
\caption{Endpoints principaux de l'API POSTIE.}
\label{tab:endpoints}
\begin{tabularx}{\linewidth}{llX}
\toprule
\textbf{Méthode} & \textbf{Endpoint} & \textbf{Rôle}\\
\midrule
POST & /api/chat & RAG conversationnel\\
POST & /api/prioritization & Scoring multi-algorithmes\\
POST & /api/prioritization/sprint-plan & Planification de sprint\\
POST & /api/agents/groom & Décomposition d'un épic en stories\\
POST & /api/documents/upload & Ingestion documentaire\\
GET  & /api/sessions/\{id\} & Historique conversationnel\\
GET  & /health & Sonde de vivacité\\
\bottomrule
\end{tabularx}
\end{table}

\section{Moteur RAG}
La chaîne RAG repose sur LangChain avec mémoire conversationnelle
glissante. Les paramètres clés (\cref{lst:rag-params}) sont issus
d'expérimentations comparatives menées en sprint~1.

\begin{lstlisting}[caption={Paramètres du moteur RAG.},label={lst:rag-params}]
chunk_size  = 800   # tokens par chunk
overlap     = 120   # chevauchement entre chunks
retrieval   = "mmr" # Maximal Marginal Relevance
k           = 6     # chunks retenus
fetch_k     = 20    # chunks pre-filtres
mem_window  = 10    # tours conversationnels conserves
\end{lstlisting}

\section{Routage multi-modèles}
La fonction \texttt{route\_model()} (\cref{lst:routage}) oriente
chaque requête vers GPT-4o (raisonnement, génération longue) ou
Mistral Large (scoring structuré, classification). Le routage est
basé sur la détection de mots-clés métiers dans la requête.

\begin{lstlisting}[caption={Routage multi-modèles simplifié.},label={lst:routage}]
SCORING_KW = ["score", "priorit", "rice", "wsjf",
              "moscow", "rank", "classif", "evalue", "noter"]

def route_model(query: str) -> ModelName:
    q = query.lower()
    if any(kw in q for kw in SCORING_KW):
        return "mistral-large"
    return "gpt-4o"
\end{lstlisting}

\section{Moteur de priorisation}
Chaque algorithme produit un score normalisé sur $[0;100]$. Le
module expose une fonction unique \texttt{prioritize()} qui
dispatche sur la fonction de scoring choisie et retourne la liste
triée. La normalisation commune permet des comparaisons
inter-algorithmes (cf. chapitre~\ref{ch:eval}).

\section{Sprint Planner}
L'algorithme suit une stratégie gloutonne sous contrainte
(\cref{lst:sprint})~:
\begin{enumerate}
  \item Trier les features par score décroissant selon
        l'algorithme choisi.
  \item Sélectionner tant que la somme cumulée d'effort
        $\leq$ vélocité.
  \item Renvoyer les listes \texttt{selected} et \texttt{deferred}
        avec le taux d'utilisation calculé.
\end{enumerate}

\begin{lstlisting}[caption={Algorithme glouton du Sprint Planner.},label={lst:sprint}]
def plan_sprint(features, velocity, algorithm):
    scored = prioritize(features, algorithm)
    selected, deferred = [], []
    cumulative = 0.0
    for rank, item in enumerate(scored, start=1):
        if cumulative + item.effort <= velocity:
            selected.append({"rank": rank, **item.dict()})
            cumulative += item.effort
        else:
            deferred.append({"rank": rank, **item.dict()})
    utilization = round(100 * cumulative / velocity, 1)
    return {"selected": selected, "deferred": deferred,
            "utilization": utilization, "velocity": velocity}
\end{lstlisting}

\section{Frontend Next.js}
L'application est organisée modulaire ment par domaine~:
\texttt{components/chat/}, \texttt{components/priority/},
\texttt{components/documents/}, \texttt{components/retro/},
\texttt{components/layout/}. L'état applicatif est géré uniquement
via les hooks React natifs (\texttt{useState}, \texttt{useReducer}),
sans librairie d'état globale, pour limiter la complexité.

\section{CI/CD}
\begin{itemize}
  \item Push sur GitHub \texttt{main} $\rightarrow$ build Vercel
        (frontend, $\approx 1$\,min).
  \item Push sur GitHub \texttt{main} $\rightarrow$ build Docker
        Render (backend, $\approx 4$\,min).
  \item Vérifications locales avant push~: \texttt{npm run build}
        et \texttt{pytest}.
\end{itemize}
